\documentclass[UTF8,a4paper,11pt,fontset=fandol]{ctexart}
\usepackage[margin=22mm]{geometry}
\usepackage{amsmath,amssymb,booktabs,longtable,array,xcolor,xurl,hyperref}
\hypersetup{colorlinks=true,linkcolor=teal,urlcolor=teal,pdftitle={Phi3Lab: Complete Supported Two-Loop BPH Analysis}}
\setlength{\parskip}{0.4em}
\setlength{\emergencystretch}{3em}
\title{全部支持两圈图的 BPH 递归分析\\\large 有限积分、原始动量绝对收敛与明确解析界}
\author{Phi3Lab 3.0 软件研究记录}
\date{2026 年 10 月 19 日}
\begin{document}
\maketitle
\begin{abstract}
本稿在六维、有共同正质量的 Euclidean 单实标量 $\phi^3$ 理论中，覆盖固定外腿标签、截肢连通 1PI 的全部两圈一、二、三点图，共十张。统一采用 BPH 递归和完整零动量 Taylor 条件；本项目尚未核验的 Z 部分不作为证明前提，森林仅作显式展开及组合互校。每种非零拓扑分别给出有限参数积分、明确可积控制、原始十二维动量绝对收敛证明，以及共同 Schwinger 正则化极限的等价桥。数值程序、独立复算及现代上游工具证据与数学证明分开记录。本稿不以一张图的证明替代任意圈数的一般定理。
\end{abstract}
\paragraph{阅读顺序}
先核对统一方案与图的物理腿映射，再读无真发散子图的原始三点图，随后比较 bubble 与 triangle 插入，最后阅读两个重叠三角子图的合并减除。参数绝对积分的常数不自动等于原始动量绝对积分的常数；数值阶数差与随机积分标准误也不是严格误差包络。
\tableofcontents
\clearpage
\section{Scheme, normalization, and the complete two-loop atlas}
\label{sec:scheme}

We work with one real scalar field in massive Euclidean $\phi^3$ theory in
six dimensions. Every internal line has the same mass squared $M=m^2>0$,
and
\[
 D(r)=\frac{1}{r^2+M},\qquad
 \int_k=\int_{\mathbb R^6}\frac{d^6k}{\pi^3},\qquad
 \langle r\rangle=\sqrt{r^2+m^2}.
\]
Unless a weight is explicitly attached, $J_R(G)=\int_k\int_q RI_G$
denotes an unweighted scalar graph integral. It excludes the fixed-external-label
symmetry weight $1/S_G$, coupling factors, and conventions for the signs of
effective-action vertices. Replacing each normalized loop measure by
$d^6k/(2\pi)^6$ multiplies a two-loop result by $(4\pi)^{-6}$.
The graph atlas consists of connected amputated 1PI multigraphs, with fixed
external labels and with tadpoles retained.

\subsection{Recursive BPH and the Taylor conditions}

For a graph with $E$ external legs, cubic valence implies the superficial
degree
\[
 \omega(G)=6L-2I=6-2E.
\]
The operator $t_G^{\omega(G)}$ takes the Taylor polynomial, through total
degree $\omega(G)$, in independent physical external momentum components
at zero, holding loop variables fixed. A subgraph Taylor operator similarly
acts on all of that subgraph's independent external momenta, including
momenta which are loop variables of the surrounding graph.

In the Bogoliubov recursion, write
\begin{align*}
 C(\gamma)&=-t_\gamma^{\omega(\gamma)}\overline R I_\gamma,\\
 \overline R I_G&=I_G+
 \sum_{\substack{\mathcal S\ne\varnothing\\
                 \mathcal S\text{ a disjoint proper spinney}}}
 I_{G/\mathcal S}\prod_{\gamma\in\mathcal S} C(\gamma),\\
 RI_G&=(1-t_G^{\omega(G)})\overline R I_G.
\end{align*}
Contractions and local polynomial insertions are understood in the second
line. Counterterms of nested subgraphs are themselves recursively prepared;
overlapping proper subgraphs cannot occur together in one disjoint spinney.
In particular, the two overlapping triangles of the two-point diamond give
two single-subgraph terms, with no product of their counterterms.
For the present atlas the overall Taylor orders are $4,2,0$ for $E=1,2,3$;
the proper bubble and triangle orders are respectively $2$ and $0$.

The terminology is BPH throughout. A primary account of the recursive
framework is A.~D.~Kennedy, \emph{A Simple Proof of the BPH Theorem}
(1996), \texttt{arXiv:hep-th/9612113}, especially its discussion of Henge
and Bogoliubov recursion in Section~3.2. A forest expansion is a useful
finite combinatorial comparison. The graph-specific parameter formulas,
explicit constants, and momentum decompositions below are derived here;
they are not attributed to Kennedy's paper. Nor does this finite atlas
constitute a new proof of an arbitrary-loop BPH theorem.

\subsection{Common regulator and the meaning of an ordinary finite integral}

Use, on every original edge, the same Schwinger lower cutoff
\[
 D_\eta(r)=\int_\eta^\infty
             e^{-\alpha(r^2+M)}\,d\alpha,\qquad \eta>0.
\]
Taylor subtraction is performed on these regulated propagators, before
Gaussian integration. Each counterterm retains the original edge labels and
the common domain $\alpha_e\ge\eta$. Thus every regulated term is integrable;
Fubini and differentiation may be used there. If
$\alpha_e=\rho x_e$, $\sum_e x_e=1$, the common radial lower limit is
\[
 T(x)=\frac{\eta}{\min_e x_e}.
\]
It is independent of external momenta. One must not replace a regulated
combination by separate unregulated integrals of its divergent terms.

For any unit direction, the absolute derivative bounds
\[
 |\partial^nD_\eta(r)|\le c_n\langle r\rangle^{-2-n},\qquad
 (c_0,c_1,c_2,c_3)=(1,2,10,72)
\]
are uniform in $\eta\ge0$. To see this, differentiate the exponential;
each term is bounded by a positive Schwinger moment. The integral from
$\eta$ to infinity is at most the corresponding integral from zero to
infinity. For example, the second derivative is bounded by
$2\langle r\rangle^{-4}+8r^2\langle r\rangle^{-6}$, and the third by
$24|r|\langle r\rangle^{-6}+48|r|^3\langle r\rangle^{-8}$.
The shift comparison
\[
 \langle r+a\rangle\ge
 \frac{\langle r\rangle}{1+|a|/m}
\]
follows from the triangle inequality applied to $(r,m)$.
These facts also bound propagator differences by the mean value formula.
They supply regulator-independent momentum majorants in the chapters below.

There are two distinct steps in each equivalence proof. A momentum-space
$L^1$ majorant proves $\int RI_\eta\to\int RI$ for the original ordinary
integral. A separate parameter-space majorant proves that the same regulated
left side converges to the displayed finite parameter formula. Identifying
these limits is legitimate; inferring original momentum $L^1$ solely from a
post-Gaussian parameter norm would not be.

\subsection{One-point graph and the finite linear condition}

For the sole two-loop one-point graph, momentum conservation fixes its only
physical external momentum to zero. There is no independent external
variable left in the prepared integrand. Consequently the overall Taylor
operator is the identity, even though its nominal order is four:
\[
 t_G^4\overline R I_{G,\eta}=\overline R I_{G,\eta},\qquad
 RI_{G,\eta}=0,
 \qquad J_R(\mathrm{L2E1G1})=0.
\]
This holds before any divergent term is integrated. The prepared graph may
contain nontrivial proper-subgraph counterterms, but the identity applies
to their complete sum. In this full zero-momentum Taylor scheme the local
linear counterterm has coefficient
$-\overline R J_{G,\eta}(0)$, with the sign translated according to the
chosen action convention. Its finite part fixes the graphwise renormalized
one-point amplitude to zero. The bare tadpole is not asserted to vanish.
This is a renormalization condition on the linear term; it is not an
independent proof of a physical vacuum expectation value condition.

\subsection{All ten fixed-label graphs and their scalar kernels}

Let $s=p^2$ in a two-point graph and let
$s_i=p_i^2$, $p_1+p_2+p_3=0$, in a three-point graph. Denote the diamond
kernel by $J_{\mathrm{ov}}(s)$, the primitive kernel by
$J_{\mathrm{pr}}(s_1,s_2,s_3)$, and the standard triangle-insertion kernel
by $J_{\mathrm{tr}}(s_1,s_2,s_3)$, where its first argument is the leg
attached to the subdivergent triangle. For a bubble insertion write
$J_{\mathrm{b},2}(s)$ or $J_{\mathrm{b},3}(a,b,c)$; in the latter,
$a,b$ are adjacent to the insertion and $c$ is opposite it.
All entries below use the same $M$.

\begin{center}
\small
\begin{tabular}{llll}
Graph & scalar kernel & proper subtraction & $S_G$\\ \hline
L2E1G1 & $0$ & bubble and triangles & $4$\\
L2E2G1 & $J_{\mathrm{ov}}(s)$ & two overlapping triangles & $2$\\
L2E2G2 & $J_{\mathrm{b},2}(s)$ & one bubble & $2$\\
L2E3G1 & $J_{\mathrm{pr}}(s_1,s_2,s_3)$ & none & $2$\\
L2E3G2 & $J_{\mathrm{tr}}(s_1,s_2,s_3)$ & triangle on leg $1$ & $1$\\
L2E3G3 & $J_{\mathrm{tr}}(s_2,s_1,s_3)$ & triangle on leg $2$ & $1$\\
L2E3G4 & $J_{\mathrm{b},3}(s_1,s_2,s_3)$ & bubble opposite leg $3$ & $2$\\
L2E3G5 & $J_{\mathrm{tr}}(s_3,s_1,s_2)$ & triangle on leg $3$ & $1$\\
L2E3G6 & $J_{\mathrm{b},3}(s_1,s_3,s_2)$ & bubble opposite leg $2$ & $2$\\
L2E3G7 & $J_{\mathrm{b},3}(s_2,s_3,s_1)$ & bubble opposite leg $1$ & $2$
\end{tabular}
\end{center}
These mappings use the actual fixed external labels. A symmetric momentum
point cannot test them: adjacent-leg exchange is a symmetry of the standard
insertion kernels, whereas exchange with their distinguished leg generally
changes the answer. The primitive kernel is symmetric under all three
external-leg permutations.

\subsection{Physical input domain and the scope of the bounds}

The mathematical three-point domain is
\[
 s_i\ge0,\qquad
 G=4s_1s_2-(s_3-s_1-s_2)^2\ge0.
\]
It is invariant under all leg permutations, and includes degenerate and
zero momenta. In the software, converted binary64 values are checked using
exact rational arithmetic. A negative requested $G$ is tolerated only when
its defect is at most $10^{-14}$ times the smallest of the three Cauchy
product scales $4ab+(c-a-b)^2$. The largest invariant is then moved inward
by binary64 steps until the effective triple has exact $G\ge0$. Requested
and effective inputs are recorded separately; all kernels and bounds use
the effective physical triple. In particular a zero leg requires equality
of the other two inputs, and an absolute error tolerance cannot hide an
inconsistent tiny leg. This numerical boundary convention does not weaken
the exact mathematical hypothesis above.

All compact-uniform statements mean that physical external momenta have a
fixed upper bound and $m\ge m_0>0$. The stronger global logarithmic scale
bounds are stated separately. Analytic amplitude bounds, original momentum
$L^1$ bounds, and quadrature refinement differences are different objects;
none is silently substituted for another. The scope is these ten massive
Euclidean graphs in the declared finite scheme, not a massless limit,
Minkowski continuation, general finite-scheme conversion, or an arbitrary-loop
induction theorem.

\clearpage
\section{The primitive two-loop three-point graph}
\label{sec:primitive}

This chapter treats L2E3G1, the graph $K_{2,3}$ with one fixed-label external
leg at the midpoint of each of its three length-two branches. Use the common
mass and normalized measures of Section~\ref{sec:scheme}. All momenta are
real Euclidean and incoming, with $p_1+p_2+p_3=0$ and $s_i=p_i^2$.
The formulas cover every such three-point configuration, including
degenerate Gram matrices and the zero point.

\subsection{Routing and subtraction}

Choose branch flows $t_1=k$, $t_2=q$, $t_3=-k-q$. The integrand is
\[
 I(p_1,p_2)=D(k)D(k+p_1)D(q)D(q+p_2)
             D(k+q)D(k+q-p_3).
\]
Every proper loop-containing 1PI subgraph has at least four internal edges;
its one-loop degree is at most $6-8=-2$. There are no proper UV-divergent
subgraphs, whereas $L=2$, $I=6$ gives overall degree zero. Therefore
\[
 \overline R I=I,\qquad RI=I(p_1,p_2)-I(0,0).
\]
The scalar answer below does not include the graph weight $1/S_G=1/2$.

\subsection{Gaussian minimization and the exact parameter polynomials}

Let the two parameters on branch $i$ be $a_i,b_i>0$, and put
$A_i=a_i+b_i$, $r_i=b_i/A_i$. Then
\[
 a_it_i^2+b_i(t_i+p_i)^2
 =A_i(t_i+r_ip_i)^2+A_ir_i(1-r_i)s_i.
\]
The constraint $\sum_i t_i=0$ makes the sum of the completed-square
vectors equal to $\sum_i r_ip_i$. Minimizing their weighted squared norm
gives the extra cost
\[
 \frac{|\sum_i r_ip_i|^2}{\sum_i A_i^{-1}}.
\]
Write $(A_1,A_2,A_3)=(A,B,C)$ and $(r_1,r_2,r_3)=(u,v,w)$. The first
Symanzik polynomial and the effective external scale are
\begin{align*}
 U&=AB+AC+BC,\\
 \Omega_p&=Au(1-u)s_1+Bv(1-v)s_2+Cw(1-w)s_3\\
 &\quad+\frac{ABC}{U}|(u-w)p_1+(v-w)p_2|^2,\\
 F&=U\{M(A+B+C)+\Omega_p\}.
\end{align*}
Here $\Omega_p$ denotes an external scale, not the superficial UV degree.
The final square can equally be written
\[
 s_1(u-w)^2+s_2(v-w)^2
 +(s_3-s_1-s_2)(u-w)(v-w).
\]
It is nonnegative exactly on the physical Gram domain. These expressions
are homogeneous before imposing any projective normalization. An independent
spanning-tree and two-forest construction of $U,F$ can therefore be checked
against them without using the same Gaussian derivation.

\subsection{Common cutoff and the finite projective integral}

Apply the same lower cutoff $\alpha_e\ge\eta$ to all six original edge
parameters and subtract the zero-momentum regulated graph first. Set
$\alpha_e=\rho x_e$, $\sum_e x_e=1$. The six-parameter Jacobian contributes
$\rho^5$, while the two six-dimensional Gaussian integrals contribute
$\rho^{-6}$. The remaining radial measure is $d\rho/\rho$, with the
same lower limit $T(x)=\eta/\min_e x_e$ for both terms.
Define
\[
 L_T(z)=\int_T^\infty\frac{d\rho}{\rho}
                  e^{-M\rho}(1-e^{-z\rho}).
\]
For $z\ge0$,
\[
 0\le L_T(z)\le\log(1+z/M),\qquad
 L_T(z)\longrightarrow\log(1+z/M).
\]
Within a branch, $da_i\,db_i=A_i\,dA_i\,dr_i$. On $A+B+C=1$ the full
projective measure is consequently
$ABC\,dA\,dB\,du\,dv\,dw$. Thus the candidate finite formula is
\begin{align*}
 J_{\mathrm{pr}}(s_1,s_2,s_3;M)
 =-&\int_{\substack{A,B\ge0\\A+B\le1}}dA\,dB
       \int_{[0,1]^3}du\,dv\,dw\\
 &\times\frac{ABC}{(AB+AC+BC)^3}
           \log\left(1+\frac{\Omega_p}{M}\right),
 \qquad C=1-A-B.
\end{align*}
The density is nonpositive; its parameter $L^1$ norm equals
$-J_{\mathrm{pr}}$. Its finiteness and the regulator limit follow from
the next estimates, rather than from separately integrating divergent terms.

For comparison only, in bare dimensional regularization with
$d=6-2\epsilon$, the radial Gamma factor is $\Gamma(2\epsilon)$. The
projective period computed below is $1/2$, giving the unweighted simple
pole $1/(4\epsilon)$. There is no extra factor of $1/2$ multiplying the
finite BPH logarithm, since
\[
 \Gamma(2\epsilon)
 \bigl(F^{-2\epsilon}-F_0^{-2\epsilon}\bigr)
 \longrightarrow-\log(F/F_0).
\]
The finite Taylor scheme is not replaced by the bare Laurent finite part.

\subsection{Exact period, compact control, and global scale behavior}

The trial branch currents $t_i=-p_i/2$ obey $\sum_i t_i=0$. The Gaussian
minimum cannot exceed their value. On $A+B+C=1$ this proves
\[
 0\le\Omega_p\le\frac{As_1+Bs_2+Cs_3}{4}
                    \le\frac{s_{\max}}4,
 \qquad s_{\max}=\max_i s_i.
\]
The weight $ABC/U^3$ is homogeneous of degree $-3$. Put
$A=x/(1+x+y)$, $B=y/(1+x+y)$, $C=1/(1+x+y)$. Its projective Jacobian is
$(1+x+y)^{-3}$, so its period is exactly
\begin{align*}
 \mathcal P
 &=\int_{A+B+C=1}\frac{ABC}{(AB+AC+BC)^3}\,dA\,dB\\
 &=\int_0^\infty\int_0^\infty
             \frac{xy\,dx\,dy}{(xy+x+y)^3}\\
 &=\int_0^\infty\frac{dy}{2(1+y)^2}=\frac12.
\end{align*}
The successive integrations are justified by Tonelli's theorem. They prove
integrability at every parameter face as well as fixing the normalization.
It follows that
\[
 |J_{\mathrm{pr}}|=\int|\text{parameter density}|
 \le\frac12\log\left(1+\frac{s_{\max}}{4M}\right)
 \le\frac{s_{\max}}{8M}.
\]
The same integrable weight times this constant controls $L_T$, uniformly
in $\eta$. Parameter dominated convergence therefore removes the common
cutoff. If $|p_i|\le P$ and $M\ge m_0^2>0$, the common bound is
\[
 |J_{\mathrm{pr}}|\le
 \frac12\log\left(1+\frac{P^2}{4m_0^2}\right).
\]
Unlike a fixed-point finiteness assertion, the logarithmic estimate also
controls arbitrarily large common external scales.

Along the ray $s_i\mapsto\lambda s_i$, $\Omega_p$ is linear in $\lambda$.
At $\lambda=1$, differentiation under the positive integrable weight gives
\[
 \left|\frac{d^nJ_{\mathrm{pr}}}{d\lambda^n}\right|
 \le\frac{(n-1)!}{2}
       \min\left\{\left(\frac{s_{\max}}{4M}\right)^n,1\right\},
 \qquad n\ge1.
\]
The first derivative is nonpositive and the second nonnegative. If the
shape is fixed and $s_{\max}$ itself is the scale variable, the bound is
$\frac12(n-1)!\min\{(4M)^{-n},s_{\max}^{-n}\}$ away from zero,
with the continuous small-scale bound at zero.

For independent vectors $z=(p_1,p_2)$, write $\Omega_p=z^TQz$. The trial
bound implies $Q\ge0$ and $\|Q\|\le1/2$, hence
\[
 |\nabla\Omega_p|^2\le2\Omega_p,\qquad
 \|\nabla^2\Omega_p\|\le1.
\]
The period is $1/2$, and elementary maximization of the differentiated
logarithm gives, with $P^2=s_{\max}$,
\begin{align*}
 \|\nabla_{(p_1,p_2)}J_{\mathrm{pr}}\|
 &\le\min\left\{\frac{P}{\sqrt2 M},
                         \frac1{2\sqrt{2M}}\right\},\\
 \|\nabla^2_{(p_1,p_2)}J_{\mathrm{pr}}\|
 &\le\frac3{4M}.
\end{align*}
These are analytic amplitude and derivative bounds, not discretization
error estimates.

\subsection{Independent absolute convergence in the original twelve dimensions}

Parameter absolute integrability after Gaussian integration alone does not
prove the original momentum $L^1$ statement. We now prove the latter.
Put $P=\max_i|p_i|$, $\Lambda=1+P/m$. The first directional derivative
bound is $|\partial D|\le2\langle r\rangle^{-3}$. Along the physical ray,
\[
 RI=\int_0^1\frac{d}{d\tau}I(\tau p_1,\tau p_2)\,d\tau.
\]
There are three terms, one for the shifted propagator on each branch.
All estimates below are uniform in $\tau\in[0,1]$.

\paragraph{One soft branch.}
Consider $|k|\ge2(|q|+P+m)$. Both lines on the $k$ and $k+q$ branches are
hard and comparable with $\langle k\rangle$. Derivatives on the soft branch
are bounded by
$512P\Lambda^3\langle k\rangle^{-8}\langle q\rangle^{-5}$;
the sum of the two hard-branch derivatives is bounded by
$2048P\Lambda^2\langle k\rangle^{-9}\langle q\rangle^{-4}$.
Since $\langle q\rangle\le\langle k\rangle$, both are controlled by
\[
 |RI|\le P C_h(\Lambda)
       \langle k\rangle^{-8}\langle q\rangle^{-5},\qquad
 C_h(\Lambda)=512\Lambda^3+2048\Lambda^2.
\]
The normalized angular factor in a six-dimensional radial integral is one.
Consequently
\begin{align*}
 \int_{|k|\ge2(|q|+P+m)}\langle k\rangle^{-8}\frac{d^6k}{\pi^3}
 &\le\frac1{8\langle q\rangle^2},\\
 \int_{\mathbb R^6}\langle q\rangle^{-7}\frac{d^6q}{\pi^3}
 &=\frac8{15m}.
\end{align*}
Permuting the three branch flows, or making a loop-variable linear change
with absolute determinant one, gives three such soft-branch regions.
Their union has absolute integral at most $PC_h/(5m)$, even if they overlap.
One of them is the soft $k+q$ region, so the channel $k\simeq-q$ is
explicitly included.

\paragraph{Comparable large branch flows.}
Let $K=\max(|k|,|q|,|k+q|)$ and $K_0=8(P+m)$. Outside the soft-branch
hierarchies, if $K\ge K_0$, all three base flows are at least $K/4$;
their physical shifts are at least $K/8$. Therefore
\[
 |RI|\le6P\,8^{13}K^{-13}.
\]
For a dyadic shell $K\in[K_n,2K_n]$, $K_n=2^nK_0$, the product of the
$k,q$ ball volumes is at most $(1024/9)K_n^{12}$. The shell integral is
bounded by a constant times $P/K_n$. Summing
$\sum_{n\ge0}K_n^{-1}=2/K_0$ gives the explicit bound
\[
 \frac{4096}{3}\frac{P\,8^{13}}{K_0}.
\]
This argument does not assume all six line momenta are independently large.

\paragraph{Compact remainder.}
For $K\le K_0$, positive mass bounds the derivative sum by $6Pm^{-13}$.
The product ball volume is at most $K_0^{12}/36$. Thus, for example,
\[
 \int_k\int_q |RI|
 \le\frac{PC_h(\Lambda)}{5m}
       +\frac{4096P\,8^{13}}{3K_0}
       +\frac{P K_0^{12}}{6m^{13}}<\infty.
\]
This deliberately coarse momentum norm differs from the much smaller
post-Gaussian parameter norm. Its controls are uniform when $P$ is bounded
and $m\ge m_0>0$.

The same propagator and first-derivative estimates hold for $D_\eta$, by
the positive Schwinger moment argument in Section~\ref{sec:scheme}. The
three regions therefore provide a uniform momentum majorant and prove
$\int RI_\eta\to\int RI$. Parameter dominated convergence identifies this
same left side with the logarithmic formula above. This completes the
equivalence to the ordinary absolutely convergent BPH momentum integral.

\subsection{A finite unit-cube parameterization}

Partition the branch-sum simplex according to the largest branch and the
order of the other two. On one of the six ordered sectors take
\[
 A=\frac{rt}{1+r+rt},\qquad
 B=\frac r{1+r+rt},\qquad
 C=\frac1{1+r+rt},\qquad 0\le r,t\le1.
\]
The exact measure transformation is
\[
 \frac{ABC}{U^3}\,dA\,dB
       =\frac{t}{(1+t+rt)^3}\,dr\,dt.
\]
Together with $u,v,w\in[0,1]$, this gives six finite five-dimensional
unit-cube integrals. The remaining reduced denominator is at least one,
and the logarithm is bounded by the already proved physical scale bound.
The six invariant permutations must be summed; equal permutations may be
combined by multiplicity. Any subsequent numerical refinement diagnostic
remains distinct from the analytic claims established here.

\clearpage
\section{Nested bubble insertions}
\label{sec:nested-bubble}

This chapter covers L2E2G2 and L2E3G4, L2E3G6, L2E3G7, with the fixed-leg
maps of Section~\ref{sec:scheme}. The internal bubble is subtracted through
second order in its external momentum. Integrating that already subtracted
subgraph does not change its finite Taylor conditions. In particular, its
finite kinetic Taylor term must be retained.

\subsection{The finite one-loop bubble}

Choose the inserted-edge momentum $k$ without physical external-momentum
dependence, and let $q$ be the bubble loop variable. Put $r=k^2$. Then
\begin{align*}
 B_R(r)&=\int_q(1-t_k^2)[D(q)D(q+k)]\\
 &=\int_0^1\left\{(M+hr)\log\left(1+\frac{hr}{M}\right)-hr\right\}\,dx,
 \qquad h=x(1-x).
\end{align*}
This formula follows with a common two-edge Schwinger cutoff: write the
parameters as $\rho(1-x),\rho x$. Their Jacobian is $\rho$ and the
Gaussian factor is $\rho^{-3}$. The Taylor-subtracted radial kernel is
\[
 \int_{T(x)}^\infty\frac{d\rho}{\rho^2}e^{-M\rho}
       (e^{-hr\rho}-1+hr\rho),\qquad
 T(x)=\frac{\eta}{\min(x,1-x)}.
\]
It is nonnegative and is controlled by its zero-cutoff value
$K(hr)$, where
\[
 K(z)=(M+z)\log(1+z/M)-z.
\]
The parameter regulator limit is therefore legitimate. Independent momentum
absolute convergence and its relation to this formula are proved below.

The finite conditions and second derivative are
\[
 B_R(0)=B_R'(0)=0,\qquad
 B_R''(r)=\int_0^1\frac{h^2}{M+hr}\,dx.
\]
These conditions uniquely fix $B_R$. In particular the term
$-r\int_0^1h\,dx=-r/6$ is part of the original finite BPH subtraction;
discarding it would change the two-loop answer. An equivalent closed form is
\[
 B_R(r)=\frac{(r+4M)^{3/2}}{3\sqrt r}
            \operatorname{arsinh}\sqrt{\frac r{4M}}
                -\frac{4M}{3}-\frac{4r}{9},
\]
with the continuous value at zero. For $|r|<4M$ the convergent expansion is
\[
 \frac{B_R(r)}M=
 \sum_{n\ge2}\frac{(-1)^n(r/M)^n(n!)^2}
                         {n(n-1)(2n+1)!},
\]
whose first term is $(r/M)^2/60$. Since
$0\le K(z)\le z\log(1+z/M)$, $h\le1/4$, and $\int_0^1h\,dx=1/6$,
\[
 0\le B_R(r)\le\frac r6\log\left(1+\frac r{4M}\right)
                    \le\frac r6\log\left(1+\frac rM\right).
\]

\subsection{A uniform original-momentum subgraph norm}

Write $K=|k|$. The explicit bound required for Fubini at two loops is
\[
 \int_q\left|(1-t_k^2)[D(q)D(q+k)]\right|
 \le K^2\left\{218+5\log\left(1+\frac{4K}{m}\right)\right\}.
\]
We give its derivation, rather than relying only on power counting.
For $K\le m$, the third-order Taylor integral remainder, the derivative
constant $c_3=72$, and
$\langle q+tk\rangle\ge\langle q\rangle/2$ imply the pointwise bound
\[
 \left|(1-t_k^2)[D(q)D(q+k)]\right|
       \le384K^3\langle q\rangle^{-7}.
\]
The factor $384$ is $(72/6)2^5$. Since
\[
 \int_q\langle q\rangle^{-7}=\frac8{15m},
\]
the norm is at most $(1024/5)K^3/m\le(1024/5)K^2<218K^2$.

For $K\ge m$, split at $R_*=2(K+m)$. On $|q|\ge R_*$ the same remainder
bound integrates to
$384K^3/R_*\le192K^2$. In the ball $|q|\le R_*$, use the following
separate integrable bounds, which are allowed on this finite region:
\begin{align*}
 \int_{|q|\le R_*}D(q)D(q+k)
 &\le\frac{R_*(R_*+K)}2\le10K^2,\\
 \int_{|q|\le R_*}D(q)^2&\le\frac{R_*^2}{2}\le8K^2,\\
 \int_{|q|\le R_*}|D(q)\,k\cdot\nabla D(q)|
 &\le2KR_*\le8K^2.
\end{align*}
The first estimate is Cauchy--Schwarz, with the shifted ball enclosed in
radius $R_*+K$. The quadratic Taylor term is bounded by
\begin{align*}
 5K^2\int_0^{R_*}\frac{Q^5\,dQ}{(Q^2+m^2)^3}
 &\le\frac52K^2\log\left(1+\frac{16K^2}{m^2}\right)\\
 &\le5K^2\log\left(1+\frac{4K}{m}\right).
\end{align*}
Here $\operatorname{area}(S^5)=\pi^3$, so the normalized radial measure
really is $Q^5dQ$. Adding $192+10+8+8$ proves the constant $218$.

\subsection{Original twelve-dimensional convergence and the regulator bridge}

At fixed $k$, the subgraph remainder is independent of physical external
momenta. The whole-graph Taylor operator therefore acts only on the shifted
outer propagators. For the two-point insertion the remaining factor is
\[
 E_2(k,p)=D(k)^2(1-t_p^2)D(k+p).
\]
If $K\ge2|p|$, the third-order remainder gives
\[
 |E_2(k,p)|\le384|p|^3\langle k\rangle^{-9}.
\]
For the three-point insertion, with adjacent momenta $P,Q$, the outer
factor is
\[
 E_3(k;P,Q)=D(k)^2\{D(k+P)D(k+Q)-D(k)^2\}.
\]
If $K\ge2\max(|P|,|Q|)$, a first-order physical-ray remainder gives
\[
 |E_3(k;P,Q)|\le64(|P|+|Q|)\langle k\rangle^{-9}.
\]
Multiplying either bound by the subgraph $q$-norm just proved, and using
the outer radial measure $K^5dK$, bounds the tail by a constant times
$K^{-2}\log K$. It is integrable. On every finite $K$ interval, positive
mass bounds the denominators and the displayed subgraph norm is finite.
Thus the complete original twelve-dimensional recursive expression is
absolutely integrable.

All these controls also hold for $D_\eta$, including the Taylor remainder
and the finite-ball bounds. The positive Schwinger moment estimates in
Section~\ref{sec:scheme} give the same constants uniformly in $\eta$.
At positive cutoff, Fubini permits the $q$ integration first. Its limit is
$B_R(k^2)$ by the one-loop parameter formula above; the two-loop majorant
permits the same limit under the remaining integral. Equivalently,
momentum dominated convergence gives the original ordinary $RI$ integral,
and the subgraph parameter limit gives the reduced formula. They coincide.
No separate unregulated divergent bare or counterterm integral is used.
The argument is uniform on bounded external momenta and $m\ge m_0>0$,
by using $m_0$ in every majorant from the outset.

This factorization is routing-sensitive. If one changes to
$k'=k+\alpha p$, the overall Taylor derivatives must also act on
$B_R((k'-\alpha p)^2)$ and on every other resulting external-momentum
dependence. Substituting $B_R(k'^2)$ instead loses finite local Taylor
terms. Similarly a subgraph loop shift applies to the complete subtracted
combination. The catalog routings can be put in the present form by
external-momentum-independent loop changes with absolute determinant one.

\subsection{Two-point insertion: positive one-dimensional integral}

Put $s=p^2$, $r=k^2$, and
\[
 R=r+M,\qquad A=R+s,\qquad W=\sqrt{(R-s)^2+4Ms}.
\]
Normalized six-dimensional angular averaging uses the density
$\frac8{3\pi}(1-z^2)^{3/2}dz$ on $-1\le z\le1$. Direct integration gives
\[
 \langle D(k+p)\rangle_{S^5}
        =\frac{4(A+2W)}{3(A+W)^2}.
\]
The odd Taylor term averages to zero, and
$\langle(k\cdot p)^2\rangle=rs/6$. Define
\begin{align*}
 F_2(r,s)&=\langle D(k+p)\rangle_{S^5}-\frac1R
                    -s\left(-\frac1{R^2}+\frac{2r}{3R^3}\right),\\
 J_{\mathrm{b},2}(s;M)
       &=\frac12\int_0^\infty\frac{r^2B_R(r)}{R^2}F_2(r,s)\,dr.
\end{align*}
The radial Jacobian is $K^5dK=\frac12r^2dr$.

For a manifestly stable sign, set $\mu=M/R$, $\sigma=s/R$, and
\[
 d=\frac{\sigma-1+\sqrt{(1-\sigma)^2+4\mu\sigma}}2,
 \qquad 0\le d\le\sigma.
\]
Algebra using $d^2+(1-\sigma)d-\mu\sigma=0$ yields
\[
 F_2(r,s)=\frac{d^2\{(1+2\mu)d^2+6\mu d+3\mu\}}
                     {3R(1+d)^2(d+\mu)}\ge0.
\]
Moreover,
\begin{align*}
 &3(1+d)^2(d+\mu)-\{(1+2\mu)d^2+6\mu d+3\mu\}\\
 &\hspace{25mm}=3d^3+(5+\mu)d^2+3d\ge0.
\end{align*}
Consequently $F_2\le d^2/R\le s^2/R^3$. The exact radial constant is
\[
 \int_0^\infty\frac{z^3\log(1+z)}{(1+z)^5}\,dz=\frac{25}{48}.
\]
For example, it is the logarithmic derivative of the beta integral
$B(4,1)=1/4$. Combining it with the bound on $B_R$ proves
\[
 0\le J_{\mathrm{b},2}(s;M)\le\frac{25s^2}{576M}.
\]
The mass dimension is two:
$J_{\mathrm{b},2}(s;M)=M J_{\mathrm{b},2}(s/M;1)$.
It vanishes with its first $s$ derivative at zero, as required by the
overall second-order momentum subtraction.

A stronger large-scale estimate uses
\[
 \frac1R-\langle D(k+p)\rangle_{S^5}
   =\frac{\sigma+2d+3d^2}{3R(1+d)^2}\ge0,
 \qquad F_2\le\frac{s}{R^2}.
\]
Split the radial integral at $r=s$. On the lower part use $F_2\le s/R^2$,
and on the upper part use $F_2\le s^2/R^3$. With
$r^3/R^4\le1/R$, $r^3/R^5\le1/R^2$, and
$L=\log(1+s/M)$, direct integration gives
\[
 J_{\mathrm{b},2}(s;M)\le
 \min\left\{\frac{25s^2}{576M},
       \frac{sL^2}{24}+\frac{s^2(L+1)}{12(M+s)}\right\}.
\]
This is at most of order $s\log^2s$ at a large common external scale.

\subsection{Three-point insertion: exact angular reduction and its bound}

Let $a=P^2$, $b=Q^2$, $c=(P-Q)^2$; $a,b$ are adjacent to the inserted
edge and $c$ is opposite it. The three fixed-label graphs use
$(P,Q)=(p_1,-p_2),(p_1,-p_3),(p_2,-p_3)$, respectively. Their physical
domain is $a,b,c\ge0$ and $ab-(a+b-c)^2/4\ge0$.
The ordinary finite integral is
\[
 J_{\mathrm{b},3}(a,b,c;M)
 =\int_k D(k)^2\{D(k+P)D(k+Q)-D(k)^2\}B_R(k^2).
\]
Use a Feynman parameter $u\in[0,1]$ and put
\[
 S=ua+(1-u)b,\qquad T=u(1-u)c.
\]
The parameter Jacobian is one and
\[
 D(k+P)D(k+Q)
 =\int_0^1\frac{du}{\{(k+uP+(1-u)Q)^2+M+T\}^2}.
\]
For $r=k^2$, let $R=r+M$, $A=R+S$, and
\[
 W=\sqrt{(R-S)^2+4MS+4rT}.
\]
The exact angular identity
\[
 \left\langle\frac1{(A+B\cos\theta)^2}\right\rangle_{S^5}
   =\frac4{(A+\sqrt{A^2-B^2})^2},\qquad
 B^2=4r(S-T),
\]
can also be obtained by differentiating the preceding first-power angular
integral at fixed $B$. It gives the two-dimensional finite representation
\[
 J_{\mathrm{b},3}
 =\frac12\int_0^\infty\frac{r^2B_R(r)}{R^2}
       \int_0^1\left\{\frac4{(A+W)^2}-\frac1{R^2}\right\}\,du\,dr.
\]
The Gram condition gives $S-T=|uP+(1-u)Q|^2\ge0$, and $0\le T\le S$.
Thus $2R\le A+W\le2(R+S)$, whence
\[
 0\le\frac1{R^2}-\frac4{(A+W)^2}\le\frac{2S}{R^3}.
\]
Using the same radial constant $25/48$ and
$\int_0^1S\,du=(a+b)/2$, we obtain
\[
 -\frac{25(a+b)}{576M}\le J_{\mathrm{b},3}(a,b,c;M)\le0.
\]
This kernel has mass dimension zero. Adjacent exchange $a\leftrightarrow b$
is a symmetry, whereas moving the opposite leg generally changes the value.

The nonnegative angular average also bounds the difference by $1/R^2$.
Let $H=\max(a,b)$, $L=\log(1+2H/M)$. Split at $r=2H$, use $1/R^2$
below and $2H/R^3$ above, and integrate as in the two-point case. The
result is
\[
 |J_{\mathrm{b},3}|\le
 \min\left\{\frac{25(a+b)}{576M},
        \frac{L^2}{24}+\frac{H(L+1)}{6(M+2H)}\right\}.
\]
At $H=0$, the physical Gram domain forces all momenta to vanish and the
continuous bound is zero. At large common scale it grows at most as
$\log^2$. Both insertion bounds are compact-uniform for bounded external
momenta and positive mass lower bound. Their sharp constants refer to the
fixed-sign angularly reduced densities; the original momentum $L^1$ proof
uses the separate, coarser subgraph estimate given above.

\clearpage
\section{Nested triangle insertions}
\label{sec:nested-triangle}

The three graphs L2E3G2, L2E3G3, L2E3G5 have one logarithmically divergent
proper triangle. They use the same massive Euclidean zero-momentum BPH
conditions as the rest of the atlas. The standard kernel below distinguishes
the physical leg $p_1$ attached to that triangle; the other two legs may be
exchanged. For the catalog graphs its arguments are, respectively,
$(s_1,s_2,s_3)$, $(s_2,s_1,s_3)$, and $(s_3,s_1,s_2)$.

\subsection{Standard routing and the four recursive terms}

All momenta are incoming, $p_1+p_2+p_3=0$. Set
\[
 a(k)=D(k)D(k+p_2)D(k-p_1),\qquad
 h(k,q)=D(q)D(q+k)D(q+k-p_1).
\]
The subgraph Taylor operator sets its independent external momenta
$(k,p_1)$ to zero at fixed $q$, so its counterterm is $-D(q)^3$.
The overall operator sets the physical $(p_1,p_2)$ to zero at fixed
$(k,q)$. Both orders are zero. Thus
\begin{align*}
 RI={}&a(k)\{h(k,q)-D(q)^3\}\\
 &-D(k)^3\{D(q)D(q+k)^2-D(q)^3\}.
\end{align*}
These four terms are the direct recursive BPH expression, not a change to
an MS finite prescription. In a standard fixed-label realization the
propagator-order edges are
$(1,3),(1,2),(2,4),(3,4),(0,3),(0,4)$, and external labels $1,2,3$ attach
to vertices $0,1,2$. This identifies the distinguished physical leg without
relying on a symmetric momentum example.

\subsection{Six common Schwinger labels and Gaussian integration}

At positive cutoff replace every propagator by $D_\eta$ before Taylor
subtraction. The four terms retain the same six original Schwinger labels
and common domain $\alpha_e\ge\eta$. Split the parameters into outer sum
$A$ and inner sum $C$, and later impose $A+C=1$. Let the outer fractions be
$(a,b,c)$ and the inner fractions $(u,v,w)$, with each triple nonnegative
and summing to one, in the propagator orders displayed above. Write
\begin{align*}
 H&=u(1-u),& B&=uvw,& D_*&=A+CH,\\
 U_{\rm raw}&=CD_*,& U_{\rm ct}&=CA,\\
 k_0&=cp_1-bp_2,&
 W&=ab s_2+ac s_1+bc s_3,\\
 G&=uv|k_0|^2+uw|k_0-p_1|^2+vw s_1,\\
 \delta&=\frac{C(AG+CBs_1)}{D_*}.
\end{align*}
The notation $D_*$ is a parameter denominator, not a propagator.
Completing the $q$ square first and then the $k$ square gives external
variance $AW+\delta$ for the raw graph and $AW$ for its subgraph
counterterm. In particular
\[
 F_{\rm raw}=U_{\rm raw}\{M(A+C)+AW+\delta\}.
\]
These identities hold before normalizing $A+C$, and may be verified by
independent spanning-tree and two-forest polynomials.

For the grouped parameters, the original six-edge measure is
\[
 A^2C^2\,dA\,dC\,da\,db\,du\,dv,
\]
with $c=1-a-b$, $w=1-u-v$. On the common projective simplex the remaining
factor is $A^2C^2\,dA$ times the two fraction-simplex measures. The radial
six-parameter measure is $\rho^5d\rho$, and the two six-dimensional
Gaussian factors supply $\rho^{-6}$. Consequently the overall subtraction
produces the logarithmic kernel
\[
 L(z)=\log(1+z/M).
\]
At positive cutoff it is instead
\[
 L_T(z)=\int_T^\infty\frac{d\rho}{\rho}
           e^{-M\rho}(1-e^{-z\rho}),\qquad
 T=\eta/\min_e x_e,
\]
and
\[
 0\le L_T(z)\le L(z),\qquad
 0\le L_T'(z)\le\frac1{M+z}.
\]
The same $T$ occurs in both paired terms. The candidate finite integral is
\begin{align*}
 J_{\mathrm{tr}}
 &=\int_0^1dA\int_{\Delta_o}da\,db\int_{\Delta_i}du\,dv\,\mathcal F,\\
 \mathcal F
 &=\frac{A^2}{C}
       \left\{\frac{L(AW)}{A^3}
                    -\frac{L(AW+\delta)}{D_*^3}\right\},\qquad C=1-A.
\end{align*}

\subsection{Exact pairing at the subdivergent face}

The two terms must be paired before $C\to0$ or the common regulator is
removed. Define $\ell(z)=\log(1+z)/z$, with $\ell(0)=1$. Then
\begin{align*}
 \mathcal F={}&
 \frac{H(D_*^2+AD_*+A^2)}{D_*^3}
       \frac WM\ell(AW/M)\\
 &-\frac{A^2(AG+CBs_1)}{D_*^4(M+AW)}
                       \ell\left(\frac\delta{M+AW}\right).
\end{align*}
There is no $1/C$ left in this stable expression. For estimates write it
as $P_+-P_-$, where the following two quantities are nonnegative:
\begin{align*}
 P_+&=\frac{A^2}{C}(A^{-3}-D_*^{-3})L(AW),\\
 P_-&=\frac{A^2}{CD_*^3}\{L(AW+\delta)-L(AW)\}.
\end{align*}
The identical decomposition with $L_T$ has the same nonnegative signs.
Mean-value estimates using $L_T'\le1/(M+z)$ give all the same majorants.
Thus the pairing remains uniform for the actual common cutoff, rather than
being an algebraic manipulation of separately divergent unregulated integrals.

\subsection{Four unit-cube sectors and their Jacobians}

There can also be joint faces $A,u\to0$ or $A,1-u\to0$. Split the inner
$u$ interval into two halves and parameterize them by
\[
 u=r/2\quad\hbox{or}\quad u=1-r/2,
 \qquad v=(1-u)z,\quad w=(1-u)(1-z).
\]
The inner Jacobian is $(1-u)/2$. Set
\[
 V=1-u,\qquad h_0=\frac{1-r/2}{2},\qquad
 H=rh_0,\qquad \frac14\le h_0\le\frac12.
\]
Also put $b_0=B/r$ by continuity; in the lower and upper halves respectively,
\[
 b_0=\tfrac12V^2z(1-z),\qquad
 b_0=\tfrac14ur\,z(1-z).
\]
In each half split the $(A,r)$ square into $A\le r$ and $r\le A$.
For $A\le r$, let $A=rt$; the Jacobian is $r$, and
\[
 D_*=rd,\quad d=t+(1-rt)h_0\ge\frac14,\quad
 N=tG+Cb_0s_1,\quad \delta=CN/d.
\]
Including the $rV/2$ Jacobian, the two positive chart densities are
\begin{align*}
 \widehat P_+&=\frac{rVh_0}{2}
       \frac{d^2+td+t^2}{d^3}\frac WM\ell(AW/M),\\
 \widehat P_-&=\frac{Vt^2N}{2d^4(M+AW)}
                  \ell\left(\frac{\delta}{M+AW}\right).
\end{align*}
For $r\le A$, let $r=At$; the Jacobian is $A$, and
\[
 D_*=Ad,\quad d=1+Ct h_0\ge1,\quad
 N=G+Ct b_0s_1,\quad\delta=CN/d.
\]
The corresponding positive densities are
\begin{align*}
 \widehat P_+&=\frac{AVt h_0}{2}
       \frac{d^2+d+1}{d^3}\frac WM\ell(AW/M),\\
 \widehat P_-&=\frac{VN}{2d^4(M+AW)}
                  \ell\left(\frac{\delta}{M+AW}\right).
\end{align*}
Finally, $a=x$, $b=(1-x)y$, $c=(1-x)(1-y)$ has outer Jacobian $1-x$.
The two inner halves and the two order sectors therefore give four
five-dimensional unit-cube integrals. All explicit radial poles are
cancelled, and the reduced denominators have the positive lower bounds
shown above. These are exact finite coordinate maps; no numerical accuracy
claim is needed for their derivation.

\subsection{An explicit linear parameter-norm bound}

Let $S=\max(s_1,s_2,s_3)$. The vector $k_0$ is a subconvex combination
of $p_1,-p_2$, and $k_0-p_1=-ap_1+bp_3$. Therefore their squared norms
are at most $S$, and
\[
 0\le W\le S/3,\qquad 0\le G\le S/3.
\]
Furthermore
\[
 HG-Bs_1=|Hk_0-uwp_1|^2\ge0,
 \qquad 0\le\delta\le CG.
\]
We bound $P_+$ and $P_-$ separately, which controls the absolute value of
their paired difference at every parameter point.

Since $D_*\ge\max(A,H)$, the stable positive term obeys $P_+\le3W/M$.
Using
$\int_{\Delta_o}W\,da\,db=(s_1+s_2+s_3)/24$ and
$\operatorname{vol}(\Delta_i)=1/2$, its integral is at most
$3S/(16M)$. For the negative term use
\[
 AG+CBs_1\le S\{AH+(A+Cu)vw\},\qquad
 \frac{A^3H}{D_*^4}\le1,
 \qquad (A+Cu)(1-u)\le D_*.
\]
The $AH$ contribution is at most $S/(4M)$ after integrating the two
fraction simplexes. The other contribution is at most $(S/(12M))I$,
where $1/12$ is the outer simplex volume $1/2$ times
$\int_0^1z(1-z)\,dz=1/6$, and
\[
 I=\int_0^1\int_0^1
 \frac{A^2(1-u)^2}{\{A+(1-A)u(1-u)\}^3}\,dA\,du.
\]
For $u\le1/2$, $A\le1/2$, use $D_*\ge A+u/4$ and extend the
$u$ integral to infinity. This part is at most
$\int_0^{1/2}2\,dA=1$. For $u\le1/2$, $A\ge1/2$, use $D_*\ge A$,
which gives $(7/24)\log2$. For $u\ge1/2$, put $v=1-u$. Since
\[
 D_*\ge\frac{A+v}{2},\qquad
 \frac{A^2v^2}{(A+v)^4}\le\frac1{16},
\]
the remaining integral is at most
$\int_0^1\int_0^{1/2}(A+v)/2\,dv\,dA=3/16$. Hence
\[
 I\le\frac{19}{16}+\frac7{24}\log2.
\]
Adding the three contributions gives the explicit parameter $L^1$ bound
\[
 \int|\mathcal F|\le C_N\frac SM,\qquad
 C_N=\frac{103}{192}+\frac{7\log2}{288}.
\]
The same pointwise nonnegative majorants hold for the common-cutoff paired
density. They prove parameter dominated convergence at all the individual
and joint UV faces, not merely at a chosen momentum point.

\subsection{A global logarithmic-square bound}

Set $T_0=S/(3M)$, $L=\log(1+T_0)$, and
\[
 f(T_0)=\int_0^{T_0}\frac{\log(1+t)}{t}\,dt
            \le L+\frac12L^2.
\]
For completeness the derivative of the difference on the right is
$(T_0-\log(1+T_0))/(T_0(1+T_0))\ge0$, and the difference vanishes at zero.
The first positive term is bounded by
\[
 P_+\le\frac3A\log(1+AT_0).
\]
The two simplex volumes are each $1/2$, so its integral is at most
$3f(T_0)/4\le3L/4+3L^2/8$.

For the second term use $\delta\le CS/3$ and
\[
 \int_{\Delta_i}\frac{A^2}{D_*^3}\,du\,dv
       \le\min\left(\frac1C,\frac1{2A}\right).
\]
The $1/(2A)$ bound follows immediately from $D_*\ge A$ and inner volume
$1/2$. For the $1/C$ bound, symmetrize the factor $1-u$ and use
$u(1-u)\ge u/2$ on $u\le1/2$; the resulting half-interval integral is
bounded by $\int_0^\infty A^2/(A+Cu/2)^3\,du=1/C$.
It follows that the second term integrates to at most
\[
 \frac12\int_0^1\frac{\log(1+CT_0)}C
                \min\left(\frac1C,\frac1{2A}\right)dA.
\]
Split at $A=1/2$. The first part is at most $L$, and the second at most
$f(T_0)/2$, giving $3L/2+L^2/4$. Altogether
\[
 |J_{\mathrm{tr}}|\le\int|\mathcal F|
 \le\min\left\{C_N\frac SM,\frac94L+\frac58L^2\right\}.
\]
This is compact-uniform when $S\le P^2$, $M\ge m_0^2>0$, and it has
at most logarithmic-square growth at a large common external scale. It is
a parameter amplitude bound, not a quadrature error or an original
momentum $L^1$ constant.

\subsection{Independent original-momentum absolute convergence}

Put $P=\max_i|p_i|$, $\Lambda=1+P/m$. On the physical ray
$p_1\mapsto tp_1$, $p_2\mapsto tp_2$, denote the factors by $a_t,h_t$.
The complete overall remainder is
\[
 RI=\int_0^1\{a_t'(h_t-D(q)^3)+a_t h_t'\}\,dt.
\]
The cancellation is made before taking absolute values. Use the uniform
propagator and first-derivative bounds of Section~\ref{sec:scheme}.

\paragraph{Hard inner loop.}
If $|q|\ge2(|k|+P+m)$, telescoping propagator differences gives
\begin{align*}
 |h_t-D(q)^3|&\le128(|k|+P)\langle q\rangle^{-7},\\
 |h_t'|&\le64P\langle q\rangle^{-7},\\
 |a_t|&\le\Lambda^4\langle k\rangle^{-6},\qquad
 |a_t'|\le4P\Lambda^5\langle k\rangle^{-7}.
\end{align*}
Thus
\[
 |RI|\le P(512\Lambda^6+64\Lambda^4)
                      \langle k\rangle^{-6}\langle q\rangle^{-7}.
\]
The hard $q$ tail is at most $1/(2\langle k\rangle)$, and
$\int_k\langle k\rangle^{-7}=8/(15m)$. This proves a finite norm in
the region that contains the logarithmically divergent inner triangle
before subtraction.

\paragraph{Hard outer loop.}
If $|k|\ge2(|q|+P+m)$, direct estimates on the first-order physical
remainder give
\[
 |RI|\le3072P\langle k\rangle^{-11}\langle q\rangle^{-2}
            +128P\langle k\rangle^{-7}\langle q\rangle^{-6}.
\]
The two hard tail integrals are bounded respectively by
$1/(160\langle q\rangle^5)$ and $1/(2\langle q\rangle)$.
The remaining $q$ integral in both cases is bounded by
$\int_q\langle q\rangle^{-7}=8/(15m)$.

\paragraph{Comparable loops, including a soft sum.}
On the remaining region let $K=\max(|k|,|q|)\ge K_0=8(P+m)$.
Then $|k|,|q|\ge K/4$, and the shifted outer momenta are at least
$K/8$. Retain $l=k+q$ as a possibly soft variable. Uniformly in $t$,
\begin{align*}
 |a_t'(h_t-D(q)^3)+a_t h_t'|
 \le{}&4P\,8^9K^{-9}D(l)D(l-tp_1)\\
          &+2P\,8^8K^{-8}D(l)\langle l-tp_1\rangle^{-3}
            +4P\,8^{13}K^{-13}.
\end{align*}
Integrating this inequality over $0\le t\le1$ bounds $|RI|$; the same
bounds below are uniform in $t$. On a dyadic shell $K\in[K_n,2K_n]$,
the $k$ ball volume is at most $64K_n^6/6$, and $|l|\le4K_n$.
Shift comparison and radial integration give
\begin{align*}
 \int_{|l|\le4K_n}D(l)D(l-tp_1)\frac{d^6l}{\pi^3}
        &\le8\Lambda^2K_n^2,\\
 \int_{|l|\le4K_n}D(l)\langle l-tp_1\rangle^{-3}
                                  \frac{d^6l}{\pi^3}
        &\le4\Lambda^3K_n.
\end{align*}
The purely hard last term uses product ball volume at most
$4096K_n^{12}/36$. Hence each shell is bounded by $P/K_n$ times the
explicit constant
\[
 \frac{1024}{3}8^9\Lambda^2
       +\frac{256}{3}8^8\Lambda^3+\frac{4096}{9}8^{13}.
\]
Since $\sum_nK_n^{-1}=2/K_0$, this region is absolutely integrable,
including $k\simeq-q$ where the two inner shifted lines are soft.

\paragraph{Compact remainder and uniform cutoff removal.}
For $K\le K_0$, positive mass gives $|RI|\le10Pm^{-13}$ and product
ball volume at most $K_0^{12}/36$. Combining the compact bound with the
three large-momentum regions proves absolute convergence in the original
twelve dimensions. This is an independent proof, not an inference from the
parameter integral.

The telescoping difference estimates and all first derivatives remain valid
for $D_\eta$, because its positive Schwinger moments are bounded by the
zero-cutoff moments. Thus these momentum majorants are uniform in $\eta$.
Momentum dominated convergence gives $\int RI_\eta\to\int RI$, while
the common-cutoff parameter pairing and the integrable majorants above give
the same left side tending to $J_{\mathrm{tr}}$. The two limits coincide.
If external momenta are bounded and $m\ge m_0>0$, use $m_0$ in the
majorants from the start to obtain compact-uniform control. The proven
result is this family of three massive Euclidean graphs, with its fixed
finite BPH scheme and physical leg mappings.

\subsection{Analytic outer-fraction reduction and endpoint smoothing}
\label{sec:triangle-reduction}

The following coordinate refinements evaluate the same paired recursive BPH
kernel. In this subsection scale $s_i$ by $M$ and set $M=1$. Retain
$C=1-A$, $H=u(1-u)$, $D_*=A+CH$, and $B=uvw$. Fix the outer fraction
$a$, put $d_o=1-a$, $b=d_oy$, $c=d_o(1-y)$. The outer Jacobian is
$d_o$. Both $Q_c=1+AW$ and $Q_r=1+AW+\delta$ are positive concave
quadratics in $y$. Writing $Q=1+v_0+q_1y+q_2y^2$, their coefficients are
\begin{align*}
 v_{0c}&=Aa d_o s_1,& q_{2c}&=-Ad_o^2s_3,\\
 q_{1c}&=A\{a d_o(s_2-s_1)+d_o^2s_3\},&&\\
 g_0&=s_1\{uvd_o^2+uwa^2+vw\},&
 g_1&=(s_1+s_3-s_2)u d_o(wa-vd_o),\\
 v_{0r}&=v_{0c}+C(Ag_0+CBs_1)/D_*,&&\\
 q_{1r}&=q_{1c}+CAg_1/D_*,& q_{2r}&=-A^2d_o^2s_3/D_*.
\end{align*}
These identities follow by expanding the already proved Gaussian variances.
In resolved sectors the factors of $D_*$ are cancelled using $D_*=rd$ or
$D_*=Ad$ before numerical evaluation.

For $b_1=q_1/(1+v_0)$, $b_2=q_2/(1+v_0)\le0$, factor
$Q/(1+v_0)=(1+z_+y)(1+z_-y)$. Then
\begin{gather*}
 z_++z_-=b_1,\quad z_+z_-=b_2,\quad
 z_+\ge0,\quad -1<z_-\le0,\\
 \mathcal L(Q)=\int_0^1\log Q(y)\,dy
   =\log(1+v_0)+f(z_+)+f(z_-),\\
 f(z)=(1+1/z)\log(1+z)-1,\quad f(0)=0,\quad f(-1)=-1.
\end{gather*}
Compute $\Delta=\sqrt{b_1^2-4b_2}$ with a hypot formula. For $b_1\ge0$
form $z_+=(b_1+\Delta)/2$ and recover $z_-=b_2/z_+$; for $b_1<0$
form $z_-=(b_1-\Delta)/2$ and recover $z_+=b_2/z_-$. Zero roots use
their continuous limits. Near zero,
$f(z)=z/2-z^2/6+z^3/12-z^4/20+\cdots$ avoids cancellation.

In either resolved sector let $\alpha,\beta$ denote the prefactors in
the preceding chart formulas, so that its paired density before the outer
Jacobian is
\[
 \alpha\frac{\log Q_c}{A}
 -\beta\frac{N}{Q_c}\ell\left(\frac{CN}{dQ_c}\right),
 \qquad \delta=CN/d.
\]
Here $N,d$ are the chart quantities already defined. Exact $y$ integration
replaces the two factors by $\mathcal L(Q_c)/A$ and
$\frac dC[\mathcal L(Q_r)-\mathcal L(Q_c)]$. If the latter logarithms
nearly cancel, use their identical positive representation
\[
 \frac dC\{\mathcal L(Q_r)-\mathcal L(Q_c)\}
 =\int_0^1\frac{N}{1+AW}
       \ell\left(\frac{CN}{d(1+AW)}\right)dy.
\]
Thus the fallback retains the finite paired kernel and removes $C$ before
evaluation. The generic reduction has four remaining integration variables.

For a large scale $S=\max s_i$, logarithmic nodes may be assigned by
\[
 A(\xi)=\frac{e^{\lambda\xi}-1}{e^\lambda-1},\quad
 A'(\xi)=\frac{\lambda e^{\lambda\xi}}{e^\lambda-1},\quad
 \lambda=\log(1+S).
\]
On the $A\le r$ sector, $r=A^{1-\eta}$, $t=A^\eta$ has
$|\partial(r,t)/\partial(A,\eta)|=-\log A$; the other sector retains
$r=At$. Smoothing $x=\sin^2(\pi\xi/2)$ has Jacobian
$(\pi/2)\sin(\pi\xi)$ and can be applied jointly to the remaining outer
fraction, the logarithmic $A$ coordinate, and the inner fraction $z$.
The upper inner half uses the exact complement $1-u=r/2$ directly.

At the physical shape $s_3=0$, $s_1=s_2=s$, one has $p_2=-p_1$ and
\[
 W=sa(1-a),\qquad
 G=s\{uv(1-a)^2+uwa^2+vw\}.
\]
Neither depends on the outer $y$. The transformation
$(a,z)\mapsto(1-a,1-z)$ leaves the complete paired kernel invariant, so
its outer weight $1-a$ can be replaced by $1/2$ after integration. The
same quadratic-log primitive now integrates $a$ exactly, leaving a
three-dimensional representation, without discarding the inner fractions.

The production implementation uses the analytic reduction on its declared
high-scale route and keeps the five-dimensional formula elsewhere. Different
endpoint maps, observed refinement trends, and extrapolations are numerical
diagnostics, not an accuracy theorem. Executed scale and coordinate checks
are recorded in \texttt{docs/nested-triangle-precision.json}; the separate
momentum-oracle evidence is in
\texttt{results/nested\_triangle\_comparison.json}. Their source fingerprints
and statistical or refinement qualifications remain part of that evidence.

\clearpage
\section{Recursive BPH subtraction of the overlapping two-loop graph}
We consider one particular massive Euclidean six-dimensional graph, rather
than infer a general BPH theorem from a finite computation. Write
\[
 D(r)=\frac1{r^2+M},\qquad M=m^2>0,\qquad s=p^2\geq0,
 \qquad d\mu(k)=\frac{d^6k}{\pi^3}.
\]
The unweighted integral has integrand
\[
 I=D(k)D(p-k)D(q)D(p-q)D(k-q).
\]
The two proper divergent subgraphs are the overlapping triangles
\(\gamma_1=\{a,b,e\}\) and \(\gamma_2=\{c,d,e\}\), each of Taylor degree zero.
The graph itself has Taylor degree two. Bogoliubov's recursion gives
\begin{align}
 \bar R I&=I-T_1I-T_2I,\\
 T_1I&=D(k)^3D(q)D(p-q),&
 T_2I&=D(q)^3D(k)D(p-k),\\
 RI&=(1-t_G^2)\bar R I.
\end{align}
The independent external inputs of the first triangle are \((q,p)\); its
loop momentum \(k\) is held fixed by \(T_1\). The second operation is symmetric.
The two triangles cannot belong to one disjoint spinney, so there is no
\(T_1T_2\) product. Expanding the recursion produces six terms, but none is
to be integrated separately without a common regulator.

\subsection{A common regulator and the finite parameter representation}
For every original edge use
\[
 D_\eta(r)=\int_\eta^\infty e^{-\alpha(r^2+M)}\,d\alpha,
 \qquad\eta>0.
\]
Apply the subgraph and overall Taylor operations to the regulated propagators
before Gaussian integration. Counterterms retain all five original Schwinger
labels, so the integration domain \(\alpha_e\geq\eta\) is common to all terms.
In the order \(D(k),D(p-k),D(q),D(p-q),D(k-q)\), call these labels
\(a,b,c,d,e\). Gaussian integration gives
\begin{align}
 U&=(a+b)(c+d)+e(a+b+c+d),\\
 V&=ab(c+d)+cd(a+b)+e(a+c)(b+d),&w&=V/U,\\
 U_1&=(a+b+e)(c+d),&w_1&=cd/(c+d),\\
 U_2&=(a+b)(c+d+e),&w_2&=ab/(a+b).
\end{align}
Set \(\alpha_e=\rho x_e\), \(\sum x_e=1\). The five-parameter measure
contributes \(\rho^4\), while the two-loop Gaussian contributes
\(\rho^{-6}\). The overall subtraction replaces \(e^{-s\rho w}\) by
\(e^{-s\rho w}-1+s\rho w\). Its common radial lower limit is
\(T(x)=\eta/\min_e x_e\). Define
\begin{align}
 K_T(z)&=\int_T^\infty\rho^{-2}e^{-M\rho}
              (e^{-z\rho}-1+z\rho)\,d\rho,\\
 K(z)=K_0(z)&=(M+z)\log(1+z/M)-z.
\end{align}
Uniformly for all \(T\geq0\) and \(z\geq0\),
\begin{equation}\label{eq:kernel-bounds}
 0\leq K_T(z)\leq\frac{z^2}{2M},\qquad
 0\leq K_T'(z)\leq\log(1+z/M)\leq\frac zM,\qquad
 0\leq K_T''(z)\leq\frac1{M+z}.
\end{equation}
These inequalities will provide a regulator-independent integrable majorant;
the limit \(\eta\downarrow0\) will therefore follow from dominated convergence.

On the simplex write
\[
 a=Au,\quad b=A(1-u),\quad c=Bv,\quad d=B(1-v),\quad e=E,
 \qquad A+B+E=1.
\]
The Jacobian is \(AB\,dA\,dB\,du\,dv\), with \(u,v\in[0,1]\). Let
\(h(t)=t(1-t)\). Then
\begin{align}
 U&=AB+E(A+B),& U_1&=B(A+E),& U_2&=A(B+E),\\
 x&=Ah(u),& y&=Bh(v),& z&=\frac{EAB}{U}(u-v)^2,\\
 w&=x+y+z.&&&
\end{align}
The finite representation is
\begin{equation}\label{eq:finite-overlap}
 J_R(s,M)=\int_{A,B\geq0,\ A+B\leq1}\!dA\,dB
       \int_0^1du\int_0^1dv\;AB\,H,
 \quad
 H=\frac{K(sw)}{U^3}-\frac{K(sy)}{U_1^3}-\frac{K(sx)}{U_2^3}.
\end{equation}
No symmetry weight, coupling, or \((4\pi)^{-6}\) loop-measure factor is
included in \(J_R\).

\subsection{An explicit absolute majorant for every parameter UV face}
Because \(EAB/U\leq AB/(A+B)\), the weighted variance identity implies
\[
 0\leq w\leq\frac{A+B}{4}.
\]
On the half-domain \(A\leq B\),
\[
 0\leq w-y\leq A,
\]
since \(z/A\leq(u-v)^2\) and \(h(u)+(u-v)^2\leq1\).
Pair the bare contribution with the first counterterm on this half-domain.
With \(r=A+E\), \(U=U_1+AE\), so
\[
 0\leq U_1^{-3}-U^{-3}\leq3AE/U_1^4.
\]
Using the mean value theorem and \eqref{eq:kernel-bounds} gives
\begin{equation}\label{eq:majorant}
 AB|H|\leq\frac{s^2}{M}(\mathcal B_1+\mathcal B_2+\mathcal B_3),
\end{equation}
where
\[
 \mathcal B_1=\frac{A^2(A+B)}{4B^2(A+E)^3},\qquad
 \mathcal B_2=\frac{3A^2E}{32B(A+E)^4},\qquad
 \mathcal B_3=\frac{B}{32(B+E)^3}.
\]
The \(u,v\) intervals have unit volume. To integrate the remaining two
parameters, set \(A=rt\), \(E=r(1-t)\), \(B=1-r\). The upper limit of
\(t\) is \(\min(1,(1-r)/r)\). Direct antiderivatives give
\begin{align}
 \int_{A\leq B}\mathcal B_1\,dA\,dB&=\frac{\log2}{6}-\frac1{96},\\
 \int_{A\leq B}\mathcal B_2\,dA\,dB&=\frac{\log2}{16}-\frac9{256},\\
 \int_{A\leq B}\mathcal B_3\,dA\,dB&=\frac1{128}.
\end{align}
The other half-domain follows by exchanging the loops. Consequently
\begin{equation}\label{eq:quadratic-bound}
 \int AB|H|\leq C\frac{s^2}{M},\qquad
 C=\frac{11\log2}{24}-\frac{29}{384}
   =0.242171624423308\ldots.
\end{equation}
This explicitly controls the triangle UV faces, their adjacent corners and
the overall radial limit. It is not just superficial power counting.

The same majorant also controls the Gaussian-integrated Schwinger expression
before radial integration: write the bare-minus-counterterm pair as
\[
 U^{-3}[E_\rho(sw)-E_\rho(sy)]
 -(U_1^{-3}-U^{-3})E_\rho(sy),\qquad
 E_\rho(z)=e^{-\rho z}-1+\rho z.
\]
The first bracket, inverse-cube difference, and \(E_\rho\) are nonnegative.
Take the triangle inequality first, and then integrate each positive piece
radially. They produce exactly the \(K\)-differences and \(K\)-values above.
Thus absolute control does not depend on cancellation after the radial
integral has already been performed.

\subsection{Uniform external-scale and derivative bounds}
For large external momentum a purely quadratic bound in \(s\) is unnecessarily
weak. Let \(L=\log(1+s/(4M))\). Use \(K'\leq L\) and
\(K(sy)\leq sBL/4\) in the same pair. The two resulting half-domain integrals
are respectively \(1/2\) and \(1/8\), so their sum contributes at most
\(5sL/8\). The remaining counterterm obeys
\[
 \int_0^{1/2}\frac{K(sA/4)}{A^2}\,dA
 \leq\frac s4\left(\ell+\frac{\ell^2}{2}\right),\qquad
 \ell=\log(1+s/(8M))\leq L.
\]
For example, integrate by parts and use
\[
 \int_0^T\frac{\log(1+t)}t\,dt
 \leq\log(1+T)+\frac12\log^2(1+T).
\]
The derivative of the right side minus the left side is nonnegative because
\(\log(1+T)\leq T\). Reflecting both half-domains yields
\begin{equation}\label{eq:global-bound}
 |J_R|\leq\int AB|H|
 \leq\min\left\{\frac{Cs^2}{M},\quad
          s\left(\frac74L+\frac14L^2\right)\right\}.
\end{equation}
Thus the global growth is at most \(s\log^2s\). For \(|p|\leq P\) and
\(m^2\geq m_0^2>0\), replacing \(s\) by \(P^2\) and \(M\) by \(m_0^2\)
gives a compact-uniform bound.

Differentiation under the parameter integral is justified by replacing
\(K(sw)\) with \(w\log(1+sw/M)\) or \(w^2/(M+sw)\) and repeating the
same majorant argument. For \(s>0\),
\begin{align}
 |J_R'(s)|&\leq\min\left\{2Cs/M,\quad1+7L/4+L^2/4\right\},\\
 |J_R''(s)|&\leq\min\left\{2C/M,\quad(5/4+L/2)/s\right\}.
\end{align}
At \(s=0\), use the continuous bound \(2C/M\); also
\(J_R(0)=J_R'(0)=0\). In vector momentum variables,
\[
 |\nabla_p J_R|\leq4CP^3/m_0^2,\qquad
 \|\nabla_p^2J_R\|\leq12CP^2/m_0^2
 \quad (|p|\leq P,\ m^2\geq m_0^2).
\]

\subsection{A separate momentum-space absolute-convergence proof}
The parameter \(L^1\) norm need not equal the momentum \(L^1\) norm.
We therefore prove absolute convergence of the original twelve-dimensional
subtracted expression independently. Write
\(\langle r\rangle=\sqrt{r^2+m^2}\). For a unit direction,
\[
 |\partial^nD(r)|\leq c_n\langle r\rangle^{-2-n},\qquad
 (c_0,c_1,c_2,c_3)=(1,2,10,72).
\]
Taylor's integral remainder applied to \(\bar R I\) gives \(P^3/6\) times
a supremum of its third directional derivative for \(|p|\leq P\).

First consider \(|k|\geq2(|q|+P+m)\). Put
\(h=D(k)D(k-p)D(k-q)\), \(a_q=D(q)D(q-p)\), and
\(\Lambda=1+P/m\). Pair before estimating:
\[
 \bar R I=(h-D(k)^3)a_q-D(q)^3D(k)D(k-p).
\]
The mean value theorem, shifted-bracket comparison, and the derivative bounds
give
\begin{align}
 |h-D(k)^3|&\leq64(P+|q|)\langle k\rangle^{-7},\\
 |\partial^nh|&\leq(64,640,9216)_n\langle k\rangle^{-6-n},
       && n=1,2,3,\\
 |\partial^ja_q|&\leq c_j\Lambda^{2+j}\langle q\rangle^{-4-j}.
\end{align}
Leibniz's formula then yields
\[
 |RI|\leq\frac{P^3C_h(\Lambda)}6
        \langle k\rangle^{-7}\langle q\rangle^{-6},
 \quad
 C_h(\Lambda)=4608\Lambda^6+1920\Lambda^4+3840\Lambda^3
                  +9216\Lambda^2+2304.
\]
The normalized radial angular factor is one, and
\[
 \int_{|k|\geq2(|q|+P+m)}\langle k\rangle^{-7}\,d\mu(k)
 \leq\frac1{2\langle q\rangle},\qquad
 \int\langle q\rangle^{-7}\,d\mu(q)=\frac8{15m}.
\]
Both symmetric hierarchy regions together therefore have absolute integral
at most
\[\frac{4C_h(\Lambda)P^3}{45m}.\]

Outside these regions, let \(K=\max(|k|,|q|)\geq K_0=8(P+m)\).
Then both loop momenta are at least \(K/4\), and their external shifts at
least \(K/8\). After three derivatives, the bare term is bounded by
\(264\,8^{11}K^{-11}D(k-q)\), and both counterterms together by
\(144\,8^{13}K^{-13}\). This does not assume that \(k-q\) is hard.
On a dyadic shell \(K\in[K_n,2K_n]\), \(K_n=2^nK_0\),
\[
 \int_{|k-q|\leq4K_n}D(k-q)\,d\mu(q)\leq64K_n^4,
 \qquad\operatorname{vol}\{|k|\leq2K_n\}=\frac{64K_n^6}{6}.
\]
The shell integrals are thus bounded by constants times \(P^3/K_n\), whose
sum is finite. One explicit bound for this comparable region is
\[
 \frac{P^3}{3K_0}
 \left[264\,8^{11}\frac{2048}{3}
              +144\,8^{13}\frac{1024}{9}\right].
\]
Finally \(K\leq K_0\) is compact. The third derivative is bounded by
\(408m^{-13}\), so \(|RI|\leq68P^3m^{-13}\); the product-ball volume is
\(K_0^{12}/36\). These three regions prove a finite momentum \(L^1\) norm.
They also give uniform control for \(|p|\leq P\), \(m\geq m_0>0\).

The same estimates hold uniformly for \(D_\eta\). Indeed, after differentiating
its Schwinger representation, the absolute values of the positive
\(\alpha\)-moment contributions from \(\eta\) to infinity are bounded by
their integrals from zero to infinity, giving the same constants
\(1,2,10,72\). The propagator-difference estimates are uniform as well.
Momentum-space dominated convergence proves
\(\int RI_\eta\to\int RI_0\); parameter-space dominated convergence proves
that the same limit is \eqref{eq:finite-overlap}. This establishes their
equality without a termwise exchange of unregulated divergent integrals.

\subsection{Interpretation}
The bounds concern this graph with a common positive mass. The explicit constant
\(C\) controls the parameter norm and amplitude magnitude; the separate,
much rougher momentum \(L^1\) estimate has different constants. Neither is a
quadrature error estimate. At \(s=M=1\), the computed benchmark is approximately
\(J_R=0.00277755520593736\); multiplying only the graph weight \(1/S=1/2\)
halves it. Couplings, standard loop-measure factors and any effective-action
sign convention remain separate. The construction follows the recursive
framework of Kennedy's BPH treatment; it does not establish a general
arbitrary-loop theorem from this representative example.

\clearpage
\section{Software, independent numerical checks and source attribution}
\subsection{A fixed scheme, not a collection of unrelated finite parts}
The software accepts the physical invariants in fixed external-label order.
For a two-point graph the sole invariant is $s=p_E^2$. For a three-point
graph they are $(s_1,s_2,s_3)=(p_{1,E}^2,p_{2,E}^2,p_{3,E}^2)$, with
$p_1+p_2+p_3=0$. The topology registry derives the inserted physical leg or
opposite physical leg from the generated graph itself.
The proper bubble retains the finite kinetic Taylor term $-s/6$; the
overall subtraction is then performed at fixed inserted-edge momentum.
The definitions in the preceding sections determine the finite scheme.
They must not be replaced by pole subtraction alone.

Shared input validation evaluates the binary64 Gram determinant using exact
rational arithmetic. A rounding tolerance is relative to the minimum of all
three Cauchy product scales, and therefore commutes with external-leg
permutations. An accepted near-boundary negative determinant is corrected by
moving the largest invariant inward until the effective binary64 Gram
determinant is nonnegative. Every kernel and analytic bound uses this same
effective physical point; the requested values, effective values and number
of correction steps are retained in the result. The larger of the first two
physical vectors is used as a reconstruction axis to avoid amplifying a
tiny-leg rounding error.

\subsection{Reference point and normalization}
At $M=1$, the two-point invariant is $s=1$ and the three-point invariants
are $s_1=s_2=s_3=1$. The following are unweighted scalar integrals in the
measure $d^6k/\pi^3\,d^6q/\pi^3$:
\begin{center}
\begin{tabular}{lr}
\toprule Topology & $J_R$\\\midrule
One-point graph & $0$ exactly\\
Overlapping two-point graph & $+0.00277755520593736$\\
Bubble-inserted two-point graph & $+0.000240153063424584$\\
Primitive $K_{2,3}$ three-point graph & $-0.0830063381643418$\\
Triangle-inserted three-point graph & $+0.0036834315027189$\\
Bubble-inserted three-point graph & $-0.00617095361238479$\\
\bottomrule
\end{tabular}
\end{center}
The three fixed-label placements in each insertion family coincide only at
this symmetric point. Asymmetric points are used to check the physical-leg
mapping. Graph weights $1/S$, vertex factors $(-g)^{E+2}$ and the ordinary
two-loop measure conversion $(4\pi)^{-6}$ remain separate.

\subsection{Independent numerical and upstream checks}
The primitive Gaussian effective scale and the nested-triangle effective
scale are checked against an independent exact spanning-tree/two-forest
construction of the graph polynomials $U,F$. The projective period
$\mathcal P=1/2$ of the primitive graph is integrated analytically and is
also reproduced numerically. Bubble insertion is checked against direct
six-dimensional angular averages, including a nonsymmetric physical point.
The overlap integral has a separate three-chart five-dimensional rational
Taylor-remainder oracle, distinct from its two-chart four-dimensional
logarithmic core.

The nested-triangle oracle first integrates the subtracted inner triangle
and then performs an outer momentum integral using a radial coordinate,
two angular coordinates and two inner simplex coordinates. Antithetic
$k\leftrightarrow-k$ samples reduce odd angular contributions. Independent
scrambled Sobol calculations use NumPy/SciPy only in the optional
verification environment. Per-scramble estimates, seeds, sample counts and
standard errors are retained. They are statistical diagnostics rather than
strict integration-error enclosures.

The official pySecDec~1.6.6 source was built in a separate environment and
used to generate and compile the original six-edge primitive graph in
$d=6-2\epsilon$. Twelve sectors remain after upstream symmetry finding.
QMC evaluations at the zero point, symmetric point and an asymmetric point
give a pole compatible with $1/(4\epsilon)$. Subtracting the zero-point
finite Laurent coefficient gives the finite BPH value for this primitive
graph, since it has no proper divergent subgraphs. In one recorded run,
the symmetric finite difference differs from the local core by
$5.37\times10^{-7}$ and the asymmetric difference by
$4.05\times10^{-7}$, meeting the recorded five-standard-deviation and
absolute-target criteria. These
are genuine upstream numerical runs; the bare finite coefficient by itself
is not called a BPH amplitude.

The task-local native build includes two syntax compatibility corrections
for the current Apple toolchain: FORM's historical \texttt{alignof} macro
does not override the C++11 keyword, and dependent SecDecUtil member
template calls supply an empty explicit template argument list. Neither
correction changes the arithmetic algorithm. The original source archive,
patch description, installed-file hashes and generated library hash are
recorded in the build and comparison JSON files.

The external feyncop source separately checks the generated graphs,
fixed-label weights and embedded divergent subgraphs. pySecDec's Python
graph interface checks all supported $U,F$ polynomials and dimensional
normalizations. None of these finite checks replaces the analytic arguments
in the preceding sections.

\subsection{Precision and reproducibility}
The core runs with Python's standard library. Optional array acceleration
and external integrators are kept separate. Quadrature order differences,
alternate coordinate maps and stochastic standard errors are explicitly
distinguished from rigorous error bounds. Analytic amplitude bounds use
outward numeric enclosures for the effective converted inputs. Inadequate
refinement or unsupported floating-point scales must be reported, rather
than inferred to be accurate from a finite return value.

The release records exact source hashes, output hashes, independently
executed tests and the freshness of each external check. A changed source
fingerprint makes previous external evidence stale until it is rerun. The
iteration ledger records substantive changes and feedback, rather than
counting repeated passing commands as new iterations. The authoritative
current counts and raw results accompany the software in \path{results/}.

\subsection{Sources and authorship}
The recursive framework follows A.~D.~Kennedy, \emph{A Simple Proof of the
BPH Theorem}, arXiv:hep-th/9612113, especially the treatment of recursive
subtraction, Taylor remainders and scale control. The explicit constants,
coordinate maps and implementation in this document were derived during
this collaboration; they are not attributed to that paper as published
results.

Modern context includes O.~Schnetz, \emph{$\phi^3$ theory at six loops},
arXiv:2505.15485, and G.~Heinrich et al., \emph{Numerical Scattering
Amplitudes with pySecDec}, arXiv:2305.19768. The former concerns minimal
subtraction renormalization functions, which are not numerically identical
to the massive finite Taylor scheme here. Exact primary-source links,
versions and the availability or execution status of published code are
retained in \path{docs/two-loop-current-sources.json} and
\path{docs/modern-references.json}.

Software, proofs and checks were produced in this collaboration. A student
report should distinguish independently understood and reproduced work,
software assistance and externally verified statements. This document
establishes the declared finite two-loop cases, not an arbitrary-loop
renormalization theorem, a massless limit or a Minkowski continuation.

\end{document}
